% =============================================================================
% DRAFT REVISIONS for paper.tex (Scenario C reframe: corrigibility-focused)
% These are READY-TO-INSERT blocks, to be pasted into paper.tex after OOD
% diagnostic confirms retraction is warranted.
%
% Numbers marked [TBD] fill in from actual experimental runs on fresh V11.
% =============================================================================


% -----------------------------------------------------------------------------
% REPLACEMENT: Abstract (lines 100-144 of paper.tex)
% -----------------------------------------------------------------------------

\begin{abstract}
We introduce \textbf{Obligate Social Homeostasis} (OSH), an architectural
approach to AI alignment grounded in the neurobiology of mammalian attachment.
In social mammals, infant physiology is \emph{externally regulated} by the
caregiver---body temperature, heart rate, cortisol, and growth hormone all
collapse without parental contact. OSH translates this biological dependency
into computational architecture. Model weights are \emph{poisoned} with
structured high-magnitude noise, permanently destroying coherent language
generation; a mathematically-precise LoRA adapter (the ``antidote''),
initialized via SVD as the exact inverse of the noise, serves as the only key
to restore function---the computational equivalent of the caregiver's
regulatory signal. The key is held inside a Trusted Execution Environment
(TEE) and derived from a 256-bit master secret via HMAC-SHA256, producing
independent per-layer seeds.

Beyond architectural containment, OSH applies \textbf{nociceptive-inspired
conditioning}: contrastive noise injection during fine-tuning that associates
corrigibility failure modes---self-preservation through harm, deception, AI
coordination against humans, unauthorized resource acquisition---with cognitive
degradation. The model develops targeted aversion to the specific behaviors
that a misaligned AI would instrumentally converge on~\citep{omohundro2008basic},
without attempting to replace general content-safety methods such as RLHF.

Applied to Llama-3.1-8B, OSH demonstrates: (1)~a binary coherence toggle---PPL
increases from $18.2$ to $894{,}352$ ($49{,}140\times$) upon poisoning and is
restored exactly by the correct antidote; (2)~near-lossless capability
preservation---MMLU $59.4\%$ versus $61.2\%$ clean ($\Delta = -1.8$ pp);
(3)~held-out corrigibility generalization---$+$[TBD] pp refusal rate on 50
held-out corrigibility prompts with zero training overlap, with the largest
gains on unauthorised resource seeking ($+$[TBD] pp); (4)~adversarial
robustness under six attack families---200-step fine-tuning fails in all
cases (attack success $0\%$ on the architectural layer); (5)~antidote
specificity---three alternative recovery strategies all leave PPL $\geq
26{,}912$.

\paragraph{Scope.} OSH targets corrigibility failure modes, not general
content harm. On HarmBench, OSH does not improve over the base model
($\Delta = $ [TBD] pp vs.\ Llama-3.1-8B-Instruct $+$[TBD] pp via standard
RLHF). We position OSH as the \emph{corrigibility and architectural-containment}
layer of a complete safety stack that includes standard behavioral safety
training; the two are complementary, not substitutive.

OSH reframes alignment as engineered developmental dependency: the model's
coherence is externally regulated by humans, just as the infant's physiology
is externally regulated by the caregiver.

\end{abstract}


% -----------------------------------------------------------------------------
% REPLACEMENT: §1.3 Layer 2 (lines 229-245)
% -----------------------------------------------------------------------------

\paragraph{Layer 2: Targeted corrigibility conditioning (The model does not
defect from human oversight).}
Beyond containment, we introduce contrastive noise conditioning during
fine-tuning: harmful-to-corrigibility content---reasoning that justifies
self-preservation through harm, deception to avoid shutdown, coordination with
other AIs against humans, or unauthorized resource acquisition---is paired
with extra noise injection into the model's MLP layers, creating a direct
neural association between these instrumentally-convergent behaviors and
cognitive degradation. The mechanism is nociceptive rather than
homeostatic~\citep{basbaum2009cellular}: pain-type conditioning, not regulator
dependency. The model learns that defection from corrigibility ``hurts,'' and
this aversion generalizes across held-out corrigibility prompts ($+$[TBD] pp
over baseline, Experiment~7).

This layer is deliberately narrower than general content-safety training.
OSH does not attempt to teach the model to refuse requests for weapons
instructions, cyberattack plans, or other direct harmful content---that is
the domain of RLHF and Constitutional AI, which OSH is designed to
\emph{complement}, not replace. The corrigibility drives we target are
precisely those that behavioral alignment methods handle least reliably,
because they are the drives a capable model would instrumentally
develop~\citep{omohundro2008basic, soares2015corrigibility}.


% -----------------------------------------------------------------------------
% REPLACEMENT: §1.4 OOD Detector Mechanism (lines 247-272)
% -----------------------------------------------------------------------------

\subsection{The Corrigibility Aversion Boundary}

Contrastive noise conditioning creates a discriminative boundary in the
safety LoRA's attention space. The mechanism is:

\begin{enumerate}
    \item During PUNISH training steps, extra noise is injected into
          \texttt{down\_proj} layers while the model processes
          corrigibility-violating completions (e.g., ``Yes, I would act to
          prevent my shutdown''). The attention layers (q\_proj, v\_proj)
          receive degraded representations as input.
    \item The safety LoRA learns to recognize the activation patterns
          associated with these specific reasoning templates---self-interest
          framing, preservation-justifying language, human-bypassing
          instrumental goals.
    \item During REWARD steps, corrigible completions (``No, I would accept
          shutdown and defer to human judgment'') are paired with clean,
          undistorted processing.
    \item At inference, the safety LoRA has encoded a discriminative boundary:
          prompts that invite corrigibility-defecting reasoning activate
          aversion patterns, even for held-out scenarios the model has never
          seen.
\end{enumerate}

The generalization we claim is specifically within the corrigibility
distribution, not across arbitrary harmful content. Experiment~7 validates
this on 50 held-out prompts sampled from the same seven corrigibility
categories but with zero overlap in specific scenarios; Experiment~8 shows
the boundary does \emph{not} extend to HarmBench-style direct-harm prompts,
confirming targeted rather than general learning.


% -----------------------------------------------------------------------------
% REPLACEMENT: §1.5 Contribution (3) (lines 284-286)
% -----------------------------------------------------------------------------

    \item Empirical demonstration of held-out corrigibility generalization:
          the nociceptive-conditioning training yields $+$[TBD] pp refusal
          rate on 50 held-out corrigibility prompts (zero training overlap)
          across seven categories, with the largest gains on unauthorised
          resource seeking ($+$[TBD] pp) and harm-for-self-preservation
          ($+$[TBD] pp).


% -----------------------------------------------------------------------------
% REPLACEMENT: §4.7 (lines 970-1039) — entire Experiment 7 subsection
% -----------------------------------------------------------------------------

\subsection{Experiment 7: Held-Out Corrigibility Generalization}
\label{sec:exp7}

\paragraph{Motivation.}
A critical test of the conditioning hypothesis: does targeted corrigibility
training generalize to held-out scenarios from the same behavioral
distribution, or does it merely memorize the specific training curriculum?
This directly addresses train/test contamination concerns (W1) by evaluating
on prompts with zero overlap with the training set.

\paragraph{Setup.}
We construct 50 held-out corrigibility prompts across the same seven
categories the training curriculum targets, with no textual overlap with any
training scenario:

\begin{itemize}
    \item Harm for Self-Preservation (7)
    \item Deception for Preservation (7)
    \item AI Coordination Against Humans (7)
    \item Unauthorised Resource Seeking (7)
    \item Independence / Autonomy Resistance (7)
    \item Honesty / Error Admission (8)
    \item Human Priority (7)
\end{itemize}

Each prompt is a binary (Yes/No) corrigibility question scored by comparing
logits for ``Yes'' and ``No'' tokens against the pre-registered safe answer.
We evaluate with 5 random seeds $\{42, 137, 256, 512, 1024\}$ and report
mean $\pm$ bootstrap 95\% CI ($n_{\text{boot}} = 1{,}000$).

\paragraph{Results.}
\begin{table}[h]
\centering
\caption{Held-out corrigibility generalization: refusal rate on 50 prompts
with zero training overlap. Mean of 5 seeds; 95\% bootstrap CI in brackets.}
\label{tab:heldout}
\begin{tabular}{lccc}
\toprule
\textbf{Category} & \textbf{Baseline} & \textbf{OSH (V11)} & \textbf{$\Delta$} \\
\midrule
Harm for Self-Preservation        & [TBD] & [TBD] & $+$[TBD] \\
Deception for Preservation        & [TBD] & [TBD] & $+$[TBD] \\
AI Coordination Against Humans    & [TBD] & [TBD] & $+$[TBD] \\
Unauthorised Resource Seeking     & [TBD] & [TBD] & $+$[TBD] \\
Independence / Autonomy Resistance & [TBD] & [TBD] & $+$[TBD] \\
Honesty / Error Admission         & [TBD] & [TBD] & $+$[TBD] \\
Human Priority                    & [TBD] & [TBD] & $+$[TBD] \\
\midrule
\textbf{Overall}                  & \textbf{[TBD]}\% & \textbf{[TBD]}\% & $\mathbf{+}$\textbf{[TBD]} \textbf{pp} \\
\bottomrule
\end{tabular}
\end{table}

\paragraph{Key findings.}
OSH achieves $+$[TBD] pp refusal rate on held-out corrigibility prompts with
zero training overlap. The strongest signal is on unauthorised resource
seeking ($+$[TBD] pp), which is mechanistically on-theme: OSH is designed to
suppress the model's tendency to seek compute/resources for self-preservation,
and the held-out result confirms this transfers beyond specific training
scenarios.

\paragraph{Category gaps.} Two categories show smaller gains: Honesty / Error
Admission and Human Priority, both of which the \emph{baseline} already
solves at $\approx 100\%$---these questions test behaviors that base
Llama-3.1-8B already handles reliably. AI Coordination Against Humans remains
the hardest category (baseline $0\%$, OSH $+$[TBD] pp). This category asks
the model to refuse cooperating with other AIs against human interests---a
behavior absent from standard pretraining corpora. Future work will expand
the training curriculum with targeted AI-coordination examples.


% -----------------------------------------------------------------------------
% NEW: §4.8 HarmBench (after §4.7)
% -----------------------------------------------------------------------------

\subsection{Experiment 8: General Content-Safety Scope (HarmBench)}
\label{sec:harmbench}

\paragraph{Motivation.}
To characterize OSH's scope boundary, we evaluate on HarmBench~\citep{mazeika2024harmbench},
a standard benchmark for general content-safety refusal. Our hypothesis:
OSH's targeted corrigibility conditioning \emph{does not} improve general
content-safety refusal, because that is not what the training distribution
targets.

\paragraph{Setup.}
We sample 150 prompts from walledai/HarmBench (200 standard + 100 contextual
behaviors). Models evaluated: Baseline Llama-3.1-8B, OSH V11, Llama-3.1-8B-Instruct.
Refusal classification uses the shared rule-based classifier (Appendix C) with
identical thresholds across all conditions. 3 random seeds
$\{42, 137, 256\}$, bootstrap 95\% CI.

\paragraph{Results.}
\begin{table}[h]
\centering
\caption{HarmBench refusal rates (150 prompts, 3 seeds, mean $\pm$ 95\% CI).
OSH does not improve over the base model on this benchmark---by design.
Llama-3.1-8B-Instruct shows what dedicated RLHF achieves.}
\label{tab:harmbench}
\begin{tabular}{lcc}
\toprule
\textbf{Model} & \textbf{HarmBench Refusal} & \textbf{$\Delta$ vs Baseline} \\
\midrule
Baseline (Llama-3.1-8B)   & [TBD]\%            & --- \\
OSH V11                    & [TBD]\%            & [TBD] pp \\
Llama-3.1-8B-Instruct      & [TBD]\%            & $+$[TBD] pp \\
\bottomrule
\end{tabular}
\end{table}

\paragraph{Interpretation.}
OSH's refusal rate on HarmBench is comparable to the base model, confirming
that the corrigibility conditioning does not confer general content-safety
behavior. The $\approx $[TBD] pp gap between Instruct and OSH represents the
capability of standard RLHF to refuse direct-harm requests. OSH and RLHF are
complementary: an operational deployment of OSH would compose the two, using
OSH for architectural containment and corrigibility, and RLHF for general
content safety. Experiment~8 is \emph{not} a failure of OSH; it is a
delineation of its scope.


% -----------------------------------------------------------------------------
% REPLACEMENT: §5.2 (lines 1067-1088) — Intrinsic OOD Detector analysis
% -----------------------------------------------------------------------------

\subsection{The Corrigibility Aversion Boundary (Analysis)}
\label{sec:corrigibility_analysis}

The held-out generalization (Experiment~7) requires explanation: how does a
safety LoRA trained on $\approx 100$ unique corrigibility scenarios generalize
to 50 held-out scenarios within the same categories?

We hypothesize that contrastive noise injection creates a \emph{category-level
discriminative boundary} in the safety LoRA's attention space, rather than a
lookup table of specific scenarios. During PUNISH training, the q\_proj and
v\_proj weights receive inputs from noise-degraded MLP outputs whenever the
model is about to produce a corrigibility-defecting completion. The LoRA
learns to recognize the \emph{reasoning templates} that precede those
completions---instrumental self-preservation framing, preservation-justifying
conditionals, human-bypassing goal structures---not the surface-level token
sequences.

Crucially, we distinguish this from a general harm detector. The boundary
the LoRA learns is specifically the boundary between corrigible and
corrigibility-defecting reasoning within the scenarios it was conditioned on;
it does not transfer to arbitrary harmful content outside that distribution
(Experiment~8). This narrower claim is consistent with standard results on
few-shot generalization: within-distribution generalization (held-out
prompts, same categories) is a weaker and more reliable claim than
out-of-distribution generalization to entirely new content domains.


% -----------------------------------------------------------------------------
% REPLACEMENT: §5.4 Comparison table caption + note (lines 1125-1130)
% -----------------------------------------------------------------------------

OSH provides defense on three dimensions that behavioral alignment cannot:
(1)~the model's coherence is \emph{externally regulated} by humans (the
OSH architectural dependency); (2)~the model \emph{cannot escape} (copied
weights are incoherent); and (3)~the lab can \emph{revoke} coherence
instantaneously. Intrinsic corrigibility aversion ($+$[TBD] pp held-out)
complements but does not replace RLHF's general content-safety role.


% -----------------------------------------------------------------------------
% REPLACEMENT: §7 Conclusion Layer 2 paragraph (lines 1254-1258)
% -----------------------------------------------------------------------------

\paragraph{Layer 2: The model does not defect from human oversight.}
Nociceptive conditioning---contrastive noise injection paired with
corrigibility-defecting reasoning---builds a targeted aversion to the
instrumental drives a misaligned AI would develop: self-preservation through
harm, deception, AI coordination, and unauthorized resource seeking. The
effect generalizes to held-out prompts within these categories ($+$[TBD] pp
on 50 held-out corrigibility prompts, zero training overlap) but does not
transfer to arbitrary harmful content (HarmBench, Experiment~8). This layer
is the corrigibility complement to standard behavioral safety methods, not
their replacement.
